
% =====================================================================
\section{The paired state from the effective action of Cornwall, Jackiw and Tomboulis}
\label{sec:cjt}
% =====================================================================

At a finite density of fermions the attraction of Sec.~\ref{sec:pairing} makes the Fermi sea unstable
against the formation of pairs. This section constructs the paired state. It first explains the
method, the effective action for composite operators of Cornwall, Jackiw and Tomboulis (CJT)
\cite{CJT1974}; then states the action of the present model; and then derives from it the gap
equation, its solutions, the two-body bound state, and the thermodynamics of the paired state. The
paired state is obtained twice, from the gap equation and from the minimum of the effective potential,
by different algorithms, and the two results are compared.

\subsection{The method}
\label{sec:method}

The usual effective action of a quantum field theory is a functional of the expectation value
$\phi(x)$ of a field. Its stationary points are the possible states, and for constant $\phi$ it
reduces to the effective potential, whose minimum is the vacuum. It is the right tool when the
quantity that acquires an expectation value is a \emph{field}. In a superfluid of fermions it is not.
A fermion field has no expectation value; what becomes non-zero is the expectation value of a
\emph{product of two fields}, a composite operator. CJT constructed the effective action for this
situation: a functional $\Gamma[\phi,G]$ of two variables, the expectation value $\phi(x)$ and the full
propagator $G(x,y)$, the connected expectation value of the product of two fields. An expectation
value $\langle\psi\psi\rangle$ is one component of the propagator of the fermions.

We use the Euclidean form, with imaginary time, because all applications below are to equilibrium.
For one real scalar field $\Phi$ with the Euclidean action $S[\Phi]$, add a source $J(x)$ for the
field and a second source $K(x,y)$ for the product of two fields,
\begin{align}
e^{W[J,K]}=\int\!\mathcal D\Phi\;\exp\Big\{&-S[\Phi]+\int_x J\Phi\nonumber\\
&+\frac12\int_x\!\int_y\Phi(x)K(x,y)\Phi(y)\Big\}.
\label{eq:WJK}
\end{align}
The derivatives of $W$ with respect to $J$ and $K$ define $\phi$ and $G$, and the CJT action is the
Legendre transform of $W$ with respect to both sources. The physical theory has no sources, and
therefore \emph{the physical $\phi$ and $G$ are the stationary point of $\Gamma$ in both variables.}
The result of CJT is that $\Gamma$ has a simple structure,
\begin{align}
\Gamma[\phi,G]={}&S[\phi]+\frac12\Tr\ln G^{-1}+\frac12\Tr\big[D^{-1}(\phi)\,G-1\big]\nonumber\\
&+\Gamma_2[\phi,G].
\label{eq:CJT}
\end{align}
The first term is the classical action. $D^{-1}(\phi)=\delta^2S/\delta\phi\,\delta\phi$ is the inverse
propagator of small fluctuations around the classical field. The traces run over the space-time
arguments and over internal indices. The last term is minus the sum of all vacuum diagrams that are
\emph{two-particle irreducible}, that is, that do not fall into two pieces when any two of their lines
are cut; it is computed with the vertices of the shifted action and, which is the essential point,
with the \emph{full propagator $G$ on every line}. Diagrams that do fall apart would count twice what
the full propagator on the lines already contains.

Varying Eq.~(\ref{eq:CJT}) with respect to $G$ gives
\begin{equation}
G^{-1}=D^{-1}(\phi)+\Sigma[\phi,G],\qquad\Sigma\equiv2\,\frac{\delta\Gamma_2}{\delta G}.
\label{eq:SD}
\end{equation}
This is the exact Schwinger-Dyson equation for the propagator, in which the self-energy $\Sigma$ is
expressed through the full propagator itself. It is called the \emph{gap equation}, because in a
superconductor its off-diagonal part determines the energy gap. For fields that are constant in space
and time, $\Gamma$ divided by the volume of space-time is the \emph{effective potential for composite
operators}, $V(\phi,G)$. At temperature $T$ the imaginary time has the period $1/T$, and the value of
$V$ at its stationary point is the thermodynamic potential per unit volume.

For a fermion field the same construction applies with two changes that follow from the
anticommuting variables: the two trace terms change sign, and the factors $\tfrac12$ are absent
because $\psi$ and $\bar\psi$ are independent. To allow pair correlations $\langle\psi\psi\rangle$,
one combines the field and its conjugate into one doubled field $\Psi=(\psi,\psi^{\dagger T})$
\cite{Nambu1960,Gorkov1958}. Its propagator $\mathcal S$ has four blocks, and the off-diagonal blocks
are the pair correlations. The components of $\Psi$ are not independent, which restores the factors
$\tfrac12$:
\begin{equation}
\Gamma[\mathcal S]=-\frac12\Tr\ln\mathcal S^{-1}-\frac12\Tr\big[\mathcal S_0^{-1}\mathcal S-1\big]+\Gamma_2[\mathcal S].
\label{eq:CJTfermion}
\end{equation}

$\Gamma_2$ is an infinite sum, and an approximation keeps some of its diagrams. The natural ordering
is by the number of loops. The diagrams with two loops give the approximations known by older names:
the Hartree approximation for a self-interacting boson, the Hartree-Fock approximation for fermions
that exchange a boson and, with pair correlations, the theory of Bardeen, Cooper and Schrieffer
\cite{BCS1957}. Even this truncation is not perturbation theory: because the lines carry the full
propagator, the gap equation sums an infinite class of perturbative diagrams self-consistently. Any
truncation of $\Gamma_2$ guarantees that the thermodynamic quantities derived from the stationary
value of $V$ are consistent with one another, and that the conservation laws of the symmetries that
the truncation respects hold \cite{LuttingerWard1960,Baym1962}. It does not guarantee that the
propagator satisfies every identity of the exact theory; Sec.~\ref{sec:thermal} meets an example.
Appendix~\ref{app:toy} tests all of this on a model that can be solved exactly, and shows that the
two-loop truncation remains accurate to {{GEN:cjt_toy_hartree_err}} per cent at couplings where
perturbation theory is useless. All applications below use the two-loop truncation.

\subsection{The action of the model}
\label{sec:action}

The fields are the fermion field and the complex Vorton-Scalar. The variables of the CJT action are
the expectation value $\phi$ of the Vorton-Scalar, defined by $\langle\PhiVS\rangle=\phi/\sqrt2$ and
real by a choice of phase, so that the fermion mass is $\gVS\phi$; the propagators $D_s$ and
$D_\vartheta$ of the radial and of the phase mode; and the propagator $\mathcal S$ of the doubled
fermion field. The action is the sum of Eq.~(\ref{eq:CJT}) for the two boson modes and of
Eq.~(\ref{eq:CJTfermion}) for the fermions,
\begin{align}
\Gamma={}&S_{\rm cl}[\phi]
+\frac12\sum_{i=s,\vartheta}\Big\{\Tr\ln D_i^{-1}+\Tr\big[D_{0i}^{-1}(\phi)D_i-1\big]\Big\}\nonumber\\
&-\frac12\Tr\ln\mathcal S^{-1}-\frac12\Tr\big[\mathcal S_0^{-1}(\phi)\,\mathcal S-1\big]\nonumber\\
&+\Gamma_2[\phi,D_s,D_\vartheta,\mathcal S].
\label{eq:action}
\end{align}
With the vertices of Eq.~(\ref{eq:Lbroken}) there are three kinds of two-particle-irreducible diagram
with two loops. \emph{Diagram (a)} is a closed fermion line with one boson line across it, for each of
the two modes. Its derivative with respect to the fermion propagator is the self-energy of a fermion
that emits and reabsorbs a boson: the normal blocks give the exchange energy of
Sec.~\ref{sec:exchange}, the pair blocks give the pairing interaction of this section and, without a
medium, the ladder equation of Sec.~\ref{sec:vacuum}. \emph{Diagram (b)} consists of two boson loops
that touch at a four-boson vertex; it gives the bosons masses that depend on the temperature
(Sec.~\ref{sec:thermal}). \emph{Diagram (c)}, two three-boson vertices joined by three boson lines, is
omitted, as in the Hartree approximation. No diagram is needed for the mean field of the Vorton-Scalar
acting on the fermions: that effect is contained in the dependence of $S_{\rm cl}$ and of
$\mathcal S_0^{-1}(\phi)$ on $\phi$, and it appears when $\Gamma$ is varied with respect to $\phi$
(Sec.~\ref{sec:matter}). Table~\ref{tab:parts} states which variables are kept, and which further
approximations are made, in each application.

\begin{table}
\caption{\label{tab:parts}The parts of the CJT action (\ref{eq:action}) that are used in each
application, and the further approximations made there.}
\begin{ruledtabular}
\begin{tabular}{>{\raggedright\arraybackslash}p{0.13\columnwidth}>{\raggedright\arraybackslash}p{0.22\columnwidth}>{\raggedright\arraybackslash}p{0.22\columnwidth}>{\raggedright\arraybackslash}p{0.31\columnwidth}}
Section & State & Variables kept & Further approximations \\
\hline
\ref{sec:cjt}, \ref{sec:order} & paired state in a medium & pair blocks of $\mathcal S$ & static exchange; particles only; normal state unchanged \\
\ref{sec:normal} & unpaired state in a medium, $T=0$ & $\phi$; normal blocks of $\mathcal S$ & exchange energy to first order; bosons at tree level \\
\ref{sec:thermal} & no medium, temperature $T$ & $\phi$, $D_s$, $D_\vartheta$ & fermions at one loop; diagram (b) only \\
\ref{sec:vacuum} & vacuum & pair blocks of $\mathcal S$ & linearized; normal blocks unchanged \\
\end{tabular}
\end{ruledtabular}
\end{table}

\subsection{The effective potential of the paired state, and the gap equation}
\label{sec:gapeq}

Keep in Eq.~(\ref{eq:action}) the pair blocks of the fermion propagator and diagram (a), with three
further approximations, each of them standard and each stated because the results depend on them.
\emph{Static exchange:} the boson propagators are evaluated at zero energy transfer. \emph{No change of
the normal state:} the diagonal part of the self-energy is absorbed into the mass $m$ and the chemical
potential $\mu$. \emph{Particles only:} the pairing of antiparticle states is neglected, which is valid
when the gap and the temperature are small compared with $m$.

A propagator with pair blocks is then specified by a gap function $\Delta(k)$. The quasiparticles
have the energies
\begin{equation}
E_k=\sqrt{\xi_k^2+\Delta(k)^2},\qquad\xi_k=\sqrt{m^2+k^2}-\mu,
\label{eq:quasi}
\end{equation}
and the pair correlation is
\begin{equation}
\kappa(k)=\big\langle b(-\vec k)\,b(\vec k)\big\rangle=\frac{\Delta(k)}{2E_k}\tanh\frac{E_k}{2T}.
\label{eq:kappa}
\end{equation}
The chemical potential is kept at its value without pairing, $\mu=\sqrt{m^2+\kF^2}$, where the Fermi
momentum $\kF$ is related to the density by $n=\nu\kF^3/6\pi^2$ with $\nu=8$ internal states. This is
correct as long as the gap is small compared with the Fermi kinetic energy
\begin{equation}
\EF=\sqrt{m^2+\kF^2}-m,
\end{equation}
the \emph{weak-coupling} regime. The sums over the Matsubara frequencies in Eq.~(\ref{eq:action}) can
be done in closed form. For a state in which all eight internal states have the same gap, which
Sec.~\ref{sec:order} calls a unitary state, and measured from the unpaired state at the same $T$ and
$\mu$, the effective potential is
\begin{align}
V[\Delta]={}&\nu\!\int\!\!\frac{d^3k}{(2\pi)^3}\Big\{\frac{|\xi_k|-E_k}{2}-T\ln\frac{1+e^{-E_k/T}}{1+e^{-|\xi_k|/T}}\nonumber\\
&\qquad\qquad\qquad+\Delta(k)\,\kappa(k)\Big\}\nonumber\\
&-\frac\nu2\int\!\!\frac{d^3k}{(2\pi)^3}\!\int\!\!\frac{d^3k'}{(2\pi)^3}\;\kappa(k)\,W(k,k')\,\kappa(k').
\label{eq:Vpair}
\end{align}
The first two lines come from the trace terms of Eq.~(\ref{eq:CJTfermion}): the energy and the
entropy of the quasiparticles, and the term $\Delta\kappa$ that the trace of
$\mathcal S_0^{-1}\mathcal S$ contributes. The last line is $\Gamma_2$, one boson exchanged between
two pairs, with the kernel
\begin{equation}
W(k,k')=\frac12\int_{-1}^{1}\!d\cos\theta\;\gVS^2\left[\frac{F_S}{q^2+m_s^2}+\sigma_X\,\frac{F_P}{q^2}\right],
\label{eq:W}
\end{equation}
where $q^2=k^2+k'^2-2kk'\cos\theta$, the amplitudes $F_S$ and $F_P$ are those of
Eq.~(\ref{eq:FSFP}), and $\sigma_X=+1$ for the vector pairs along the extra axes and $\sigma_X=-1$ for
their two-form pairs. The angular integral is done in closed form.

Because the derivative of the first two terms in the braces with respect to $\Delta$ is $-\kappa$,
\begin{equation}
\frac{\delta V}{\delta\Delta(k)}=\nu\,\frac{\partial\kappa}{\partial\Delta}\,\Big[\Delta(k)-\int\!\!\frac{d^3k'}{(2\pi)^3}W(k,k')\,\kappa(k')\Big],
\label{eq:dV}
\end{equation}
and $\partial\kappa/\partial\Delta>0$. The potential is therefore stationary exactly where the
\emph{gap equation}
\begin{equation}
\Delta(k)=\frac{1}{4\pi^2}\int_0^\infty\!dk'\,k'^2\;W(k,k')\;\frac{\Delta(k')}{E_{k'}}\,\tanh\frac{E_{k'}}{2T}
\label{eq:gap}
\end{equation}
holds. A unit test of the program checks Eq.~(\ref{eq:dV}) against numerical derivatives of
Eq.~(\ref{eq:Vpair}).

\subsection{Solutions of the gap equation}
\label{sec:gapsol}

The gap equation was solved by iteration on a momentum grid that is refined logarithmically around
$\kF$, for {{GEN:gap_n_points}} combinations of the coupling $\gVS$, the mass ratio $m_s/m$ and the
Fermi momentum $\kF/m$, and for three kernels: the vector channel ($\sigma_X=+1$), the radial mode
alone, and the two-form channel ($\sigma_X=-1$). The critical temperature $T_c$ is the temperature at
which the linearized equation has the eigenvalue 1. The strength of each mode is measured by the
dimensionless coupling
\begin{equation}
\lambda=N(0)\,W(\kF,\kF),\qquad N(0)=\frac{\kF\sqrt{m^2+\kF^2}}{2\pi^2},
\label{eq:lambda_coupling}
\end{equation}
where $N(0)$ is the density of states at the Fermi surface per internal state; $\lambda_S$ and
$\lambda_P$ denote the contributions of the radial and of the phase mode. Two validations precede the
results. For a kernel that is constant in a shell around the Fermi surface the program reproduces the
textbook results $\Delta=2\omega_c\,e^{-1/\lambda}$ and $T_c=0.567\,\Delta$. And in the scan the
ratio of $T_c$ to the gap at the Fermi momentum lies between {{GEN:gap_tc_lo}} and {{GEN:gap_tc_hi}}
wherever the coupling is weak.

\begin{figure*}
\includegraphics[width=\textwidth]{fig_gap}
\caption{\label{fig:gap}Gap at the Fermi surface in units of the Fermi kinetic energy, as a function
of the Fermi momentum, for two sets of couplings. In every case the vector pair along an extra axis,
to which the radial mode and the phase mode both contribute attraction, has the largest gap; the
two-form pair, in which the phase mode repels, has the smallest. Open symbols mark points above the
dashed line, where the gap is no longer small compared with the Fermi energy and the gap equation with
a fixed chemical potential is not quantitative.}
\end{figure*}

\begin{table*}
\caption{\label{tab:gap}Solutions of the gap equation (\ref{eq:gap}) at zero temperature for nine of
the {{GEN:gap_n_points}} parameter sets: the dimensionless couplings $\lambda_S$ and $\lambda_P$ of
the radial and of the phase mode, the gap at the Fermi momentum in units of the Fermi kinetic energy
for the three kernels, and the ratio of the critical temperature to the gap for the vector channel.
An entry ``$<10^{-7}$'' is below the resolution of the calculation.}
\begin{ruledtabular}
\begin{tabular}{ccccccccc}
$\gVS$ & $m_s/m$ & $\kF/m$ & $\lambda_S$ & $\lambda_P$ & vector: $\Delta/\EF$ & radial only & two-form & $T_c/\Delta$ (vector) \\
\hline
{{ROWS:gap_table|cols=0-8}}
\end{tabular}
\end{ruledtabular}
\end{table*}

Figure~\ref{fig:gap} and Table~\ref{tab:gap} show the result that matters for hypothesis H1. \emph{The
gap of the vector channel is the largest} in all {{GEN:gap_n_vector_largest}} parameter sets, and it
exceeds that of the two-form channel by a factor that grows with the density and with the mass of the
radial mode, from {{GEN:gap_ratio_lo}} to {{GEN:gap_ratio_hi}} in the weak-coupling part of the scan.
The phase mode alone, that is, the neutral current without the radial mode, gives a gap that is
negligible at low density, because its coupling $\lambda_P=N(0)\,\gVS^2/4(m^2+\kF^2)$ vanishes with
$\kF$; it becomes comparable to the contribution of the radial mode when $\kF\sim m$.

\subsection{Two fermions: the bound pair}
\label{sec:dimer}

At low density the attraction (\ref{eq:yukawa}) is felt by two fermions alone. Whether they bind is
decided by one number,
\begin{equation}
\DVS\equiv\frac{\aVS\,m}{m_s}.
\label{eq:DVS}
\end{equation}
The radial Schr\"odinger equation of the relative motion, with the reduced mass $m/2$, is
$u''+\DVS\,e^{-\rho}u/\rho=\varepsilon u$ in the variables $\rho=m_sr$ and $\varepsilon=mE_b/m_s^2$,
where $E_b$ is the binding energy. Its numerical integration gives a bound state if and only if
\begin{equation}
\DVS>D_c={{GENM:dimer_dc}},
\label{eq:Dc}
\end{equation}
which is the critical value known for the Yukawa potential \cite{RogersGraboskeHarwood1970}. The
scattering length has the Born limit $a=-\DVS/m_s$ for weak coupling and diverges at $D_c$
(Fig.~\ref{fig:dimer}).

\begin{figure*}
\includegraphics[width=\textwidth]{fig_dimer}
\caption{\label{fig:dimer}Two fermions in the Yukawa potential of the radial mode. (a) Scattering
length in units of $1/m_s$; it diverges at the critical coupling $D_c$, where the bound state appears.
(b) Binding energy of the pair.}
\end{figure*}

For $\DVS<D_c$ the superfluid is of the kind described by Bardeen, Cooper and Schrieffer at all
densities: large, overlapping pairs, described by Eq.~(\ref{eq:gap}). For $\DVS>D_c$ a dilute gas
consists of bound pairs, bosons of spin 0, that condense; as the density rises the pairs overlap, and
the state crosses over continuously to the first kind \cite{NozieresSchmittRink1985}. In the scan of
Sec.~\ref{sec:gapsol} the {{GEN:gap_n_beyond}} points marked by open symbols are on the way to this
regime. There the \emph{existence} of the paired state is not in doubt, since two fermions already
bind, but the numbers obtained from Eq.~(\ref{eq:gap}) are not quantitative.

The bound state and the instability of the vacuum against pairing are the same thing. Without a Fermi
sea, at $T=0$ and $\mu<m$, the linearized gap equation becomes, for momenta small compared with $m$ and
with $\mu=m-E_b/2$, the Schr\"odinger equation of the relative motion in momentum space. A bound
state appears when the largest eigenvalue of the kernel at $\mu=m$ reaches 1. Solving this condition
with the relativistic kernel (\ref{eq:W}) gives Table~\ref{tab:threshold}: for a force of long range
compared with the Compton wavelength of the fermions the threshold is
{{GEN:cjt_threshold_best}}, against {{GEN:cjt_threshold_ode}} from the Schr\"odinger equation in
coordinate space. The two programs share no code. For $m_s$ comparable to $m$, relativistic effects
raise the threshold by a few per cent.

\begin{table}
\caption{\label{tab:threshold}Critical value of $\DVS=\aVS m/m_s$ for a bound pair, obtained from the
gap equation in vacuum with the relativistic kernel, and its relative difference from the value
{{GEN:cjt_threshold_ode}} of the non-relativistic Schr\"odinger equation.}
\begin{ruledtabular}
\begin{tabular}{ccc}
$m_s/m$ & $D_c$ from the gap equation & relative difference \\
\hline
{{ROWS:cjt_threshold_table}}
\end{tabular}
\end{ruledtabular}
\end{table}

\subsection{Independent checks through the effective potential}
\label{sec:checks}

The potential (\ref{eq:Vpair}) was evaluated on the momentum grid for twelve of the parameter sets,
eight in the vector channel and two each in the two-form channel and with the radial mode alone.
Table~\ref{tab:cjtchecks} gives four checks.

\emph{An independent solution.} The gap equation is solved by Newton's method with the exact
Jacobian, starting from a function of different height and shape, where the scan of
Sec.~\ref{sec:gapsol} used iteration. The two gap functions agree at every grid point within
{{GEN:cjt_pair_worst_newton}} of the gap, and with the stored results of the scan within
{{GEN:cjt_pair_worst_scan}}.

\emph{The solution is a minimum.} Along the family $s\,\Delta(k)$ of rescaled gap functions the
potential has its minimum at $s=1$, within {{GEN:cjt_pair_worst_smin}}. All second variations are
positive: the matrix of second derivatives is positive definite, and its smallest eigenvalue, in
suitable units, is given in the table.

\emph{The value at the minimum} is the condensation energy. In weak coupling it must approach
$V_{\rm BCS}=-\nu N(0)\Delta_F^2/4$, where $\Delta_F$ is the gap at the Fermi momentum. The ratio lies
between {{GEN:cjt_pair_v_lo}} and {{GEN:cjt_pair_v_hi}}, closest to 1 for the smallest gaps.

\emph{The critical temperature.} Below $T_c$ the unpaired state $\Delta=0$ is unstable: the second
variation of $V$ there has a negative direction. The temperature at which its determinant changes
sign agrees with the $T_c$ of the linearized gap equation within {{GEN:cjt_pair_worst_tc}}, the
resolution of the bisection.

\begin{table*}
\caption{\label{tab:cjtchecks}Checks of the paired state through the CJT potential (\ref{eq:Vpair}).
$m_s$ and $\kF$ are in units of $m$. ``Newton'' and ``scan'' are the largest relative differences
between the gap function found by Newton's method on the potential and, respectively, the gap function
found by iteration and the gap stored by the scan of Sec.~\ref{sec:gapsol}. $V/V_{\rm BCS}$ is the
condensation energy relative to its weak-coupling value. The last column is the relative difference
between the two determinations of $T_c$.}
\begin{ruledtabular}
\begin{tabular}{ccclccccccc}
$\gVS$ & $m_s$ & $\kF$ & pair & $\Delta/\EF$ & Newton & scan & $V/V_{\rm BCS}$ & smallest curvature & $T_c/\Delta$ & $T_c$: two methods \\
\hline
{{ROWS:cjt_pairing_table}}
\end{tabular}
\end{ruledtabular}
\end{table*}

\emph{The Ginzburg-Landau coefficients.} Near $T_c$ the potential can be expanded along a gap function
of fixed shape, $\Delta(k)=s\,\varphi(k)$, where $\varphi$ is the solution of the linearized gap
equation normalized to 1 at the Fermi momentum: $V=\alpha(T)\,s^2+\tfrac12\beta\,s^4+\dots$ Expanding
Eq.~(\ref{eq:Vpair}) gives
\begin{align}
\alpha(T)&=\frac\nu2\big[1-\eta(T)\big]\!\int\!\!\frac{d^3k}{(2\pi)^3}\,\kappa'\varphi^2,\nonumber\\
\frac\beta2&=\nu\Big[\frac18-\frac{\eta(T)}{6}\Big]\!\int\!\!\frac{d^3k}{(2\pi)^3}\,\kappa'''\varphi^4,
\label{eq:GLcjt}
\end{align}
where $\eta(T)$ is the largest eigenvalue of the linearized gap equation, with $\eta(T_c)=1$, and
$\kappa'$, $\kappa'''$ are the first and third derivatives of Eq.~(\ref{eq:kappa}) with respect to
$\Delta$ at $\Delta=0$. With a constant density of states these expressions reduce at $T_c$ to the
weak-coupling coefficients $a$ and $b$ of Eq.~(\ref{eq:GL}) below, with $\alpha=8a$ and
$\beta=8b$. The gradient coefficient $K$ is obtained independently, from the dependence of the pair
susceptibility on the momentum of the pair. Table~\ref{tab:gl} compares the three coefficients,
extracted numerically at $T_c$, with the weak-coupling formulas: they agree within
{{GEN:cjt_gl_worst}} per cent, and the deviations decrease with the coupling.

\begin{table}
\caption{\label{tab:gl}The Ginzburg-Landau coefficients extracted from the CJT potential at $T_c$,
divided by the weak-coupling formulas of Eq.~(\ref{eq:GL}): $a$ and $b$ from the expansion of the
potential, $K$ from the pair susceptibility.}
\begin{ruledtabular}
\begin{tabular}{cccccc}
$\gVS$ & $m_s/m$ & $\kF/m$ & $a$ & $b$ & $K$ \\
\hline
{{ROWS:cjt_gl_table|cols=0,1,2,5,6,7}}
\end{tabular}
\end{ruledtabular}
\end{table}

\subsection{Thermodynamics of the paired state}
\label{sec:thermo}

At the solution of the gap equation the potential (\ref{eq:Vpair}) is the difference of the
thermodynamic potentials of the paired and of the unpaired state. Figure~\ref{fig:thermo} and
Table~\ref{tab:thermo} show it, together with the gap, the entropy $S$ and the heat capacity
$C=T\,dS/dT$, for $\gVS=3$, $m_s=m$ and $\kF=0.5\,m$ in the vector channel.

\begin{figure*}
\includegraphics[width=\textwidth]{fig_cjt_thermo}
\caption{\label{fig:thermo}The paired state as a function of temperature, for $\gVS=3$, $m_s=m$,
$\kF=0.5\,m$. (a) Gap at the Fermi momentum. (b) Condensation energy, that is, minus the CJT
potential, in units of its value at $T=0$. (c) Heat capacity in units of that of the unpaired state
at the same temperature; it jumps at $T_c$.}
\end{figure*}

\begin{table}
\caption{\label{tab:thermo}The paired state as a function of temperature (parameters of
Fig.~\ref{fig:thermo}): gap, condensation energy, and entropy and heat capacity relative to the
unpaired state (subscripts $s$ and $n$).}
\begin{ruledtabular}
\begin{tabular}{ccccc}
$T/T_c$ & $\Delta(T)/\Delta(0)$ & $-V(T)/\lvert V(0)\rvert$ & $S_s/S_n$ & $C_s/C_n$ \\
\hline
{{ROWS:cjt_thermo_table}}
\end{tabular}
\end{ruledtabular}
\end{table}

The universal numbers of weak coupling are reproduced, with corrections of the size expected for a
gap of a few per cent of the Fermi energy:
\begin{align}
\frac{\Delta(0)}{T_c}&={{GENM:cjt_gap0_over_tc}}\quad(1.764),\nonumber\\
\frac{C_s-C_n}{C_n}\Big|_{T_c}&={{GENM:cjt_c_jump}}\quad(1.426),
\label{eq:universal}
\end{align}
where the values in parentheses are those of the weak-coupling limit. They hold for both spacetimes
and for any unitary state. Two consequences are used in Sec.~\ref{sec:dm}. The entropy of the paired
state is exponentially small below about $T_c/4$, so that the paired fluid is cold in the strict
sense. And the condensation energy per fermion, $\tfrac38\Delta^2/\EF$ in the non-relativistic limit,
is a small fraction of the Fermi energy, so that pairing changes the pressure only by terms of
relative order $(\Delta/\EF)^2$.

\subsection{What has been established, and what has not}

Established: within the approximations of Sec.~\ref{sec:gapeq} the gap equation has solutions; the
vector channel has the largest gap; the numerical method reproduces the textbook limits; two fermions
bind when $\DVS>D_c$; the solution of the gap equation is the minimum of the CJT potential, and a
different formulation with different algorithms reproduces the gap functions, the critical
temperatures, the condensation energy, the critical coupling of the bound pair and the three
Ginzburg-Landau coefficients, to the accuracies of Tables~\ref{tab:threshold} to \ref{tab:gl}.

Not established: anything beyond the two-loop truncation of $\Gamma_2$ and the three approximations
of Sec.~\ref{sec:gapeq}. In particular the dependence of the pairing interaction on the energy
transfer, which is known to change the prefactor of the gap in comparable systems; the
renormalization of the normal state; fluctuations of the order parameter; and the strong-coupling
regime beyond the open symbols of Fig.~\ref{fig:gap}.
